% SOURCE_AUTHORITY: sga5_fr_workpass.tex, lines 9694--9721 (LF-normalized unit).
% SOURCE_SHA256: 791F4EFFC5E02832D5D77ED03518C8156D6F07E4C8238B03545DB93D883FBB28
\subsubsection*{3.9}
Sejam $X$ um esquema e $E$ um $\mathcal O_X$-módulo localmente livre. Um passo imediato para o limite permite definir, para todo número primo $\ell$ primo com as características residuais de $X$, um classes de Chern $\ell$-ádicas
\[
c_i(E)\in H^{2i}(X,\mathbb Z_\ell(i))
\stackrel{\mathrm{def}}{=}\varprojlim_\nu H^{2i}(X,\mu_{\ell^\nu}^{\otimes i}),
\]
que possuem as propriedades usuais explicadas em 3.4 e 3.7 e que, no contexto de 3.8, são manifestamente compatíveis, em um sentido evidente, com as classes de Chern definidas por métodos topológicos.

Tensorizando por $\mathbb Q_\ell$, deduz-se uma torção classes de Chern $\ell$-ádicas módulo
\[
c_i(E)\in H^{2i}(X,\mathbb Q_\ell(i))
=H^{2i}(X,\mathbb Z_\ell(i))\otimes_{\mathbb Z_\ell}\mathbb Q_\ell,
\]
que verificam claramente as mesmas propriedades.

Agora vamos supor que $X$ é um esquema quase projetivo e suave sobre um corpo $k$, e vamos denotar aqui por $C^*(X)$ seu anel de Chow (Sem. Chevalley: Anneaux de Chow et applications III 4). Lembre-se que definimos em IV um morfismo de anéis graduados
\[
\varphi_X:C^{\bullet}(X)\to A^{\bullet}(X)
\]
Compatível com morfismos de imagem direta e de imagem inversa. Por outro lado, o isomorfismo canônico
\[
\varepsilon_X:\Pic(X)\xrightarrow{\sim} C^1(X)
\]
que associa a toda $\mathcal O_X$-módulo inversível $L$ a classe do divisor de uma seção racional de $L$ que não é nula sobre nenhum componente conexo de $X$, dá lugar, seguindo o modelo do artigo de GROTHENDIECK antes citado, onde se supôs que $k$ era algebraicamente fechado, a uma teoria de classes de Chern com valores em $C^*(X)$. O composto morfismo
\[
\varphi_X\circ\varepsilon_X:\Pic(X)\longrightarrow H^2(X,\mathbb Z_\ell(1))
\]
Não é mais do que a primeira classe de Chern $\ell$-ádica. Com efeito, basta verificar isso sobre os elementos de $\Pic(X)$ da forma $\check J$, com $J$ um ideal inversível de $\mathcal O_X$, e então segue-se de (3.6.4). De propriedades de unicidade das teorias de classes de Chern deduz-se que o morfismo $\varphi_X$ transforma classes de Chern em classes de Chern.
